%\TEX program = pdflatex
\documentclass[a4paper]{report}
%[UTF8]{ctexart}

\usepackage[utf8]{inputenc}       % Windows
\usepackage[english, russian]{babel}
\usepackage[T2A]{fontenc}

\usepackage{graphicx}             % Для включения графических изображений
\usepackage{fancyhdr}             % Для тонкой настройки колонтитулов
\usepackage{makeidx}              % Для предметного указателя
\usepackage{tikz}                 % Для добавления меток к скриншотам
\usetikzlibrary{calc}
\usetikzlibrary{backgrounds,fit}  % Для заднего фона
\usetikzlibrary{shadows}          % Для теней
\usetikzlibrary{arrows.meta}

\usepackage{pscyr}
\pdfmapfile{+pscyr.map}
\renewcommand{\rmdefault}{fjn}
\renewcommand{\sfdefault}{far}
\renewcommand{\ttdefault}{fcr}

% Создаем предметный указатель
\makeindex

% \title{123123}
% \author{123123}
% \date{\today}
\begin{document}

% Создаем
% \maketitle

% Окружения
\newenvironment{Titul}
{
    \begin{center}}{\end{center}
    \baselineskip 12pt
    \addtolength{\textheight}{-14mm}
}

\newenvironment{Anot}
{
    \begin{quote}\settowidth{\parindent}{~~~}\scriptsize\mbox{~~~}}{\end{quote}
}

%% title, authors, affilations, anotation
%% title
% \title{Лабораторные работы}
% \tnotetext[titlenote]{Комментарий к статье. Таких комментариев может быть несколько.}
\thispagestyle{empty}
\begin{Titul}
{
    \begin{center}
        NIIMAILTAH \\
    \end{center}
    \vspace{8em}

    \bf{Статья по использованию коммуникационных интерфейсов}
    \footnote{Статья написана с ознакомительной и справочной целью.}
    \\[4ex]
    \begin{flushright}
    Niimailtah
    \end{flushright}
}
\end{Titul}    
\vspace{4em}
\begin{Anot}
    Первоначальный черновой набор текста закончен.
    
    Ключевые слова: список слов не более 10.
    I2C, TWI, SPI, UART.
\end{Anot}

\vspace{\fill}

\begin{center}
    Челябинск \\2026
\end{center}

\pagebreak

\tableofcontents
\pagebreak

% \printglossaries
% \pagebreak

% ====== Основной текст статьи ======
\chapter{Введение}
% Введение
\section{Введение}

Для того, чтобы сделать PDF-документ требуется установить MiKTeX, VSCode с плагином LaTeX Workshop, дополнительно еще установил Perl от Strawberry.  
Далее идем настраивать LaTeX: фонт,  

% \pagebreak

\chapter{I2C}
% I2C
\section{I2C}

I2C

% \pagebreak

\chapter{SPI}
% SPI
\section{SPI}

SPI

% \pagebreak


\end{document}
